% ---
% id: meat-core/ass:separate-executions
% kind: ass
% title: "Executions that share no key and no device state"
% summary: >-
%   Several executions, fixed in advance and run against one adversary, touch
%   none of each other's registers, leave their keys unused until all have
%   ended, and keep their own premises with the others counted as the
%   adversary's. It is what lets their security parameters add.
% prominence: standard
% status: accepted
% uses: [def:execution-end]
% assumes: []
% refs:
%   - "portmann2022security, Sec. II.E"
% provenance: newly-derived
% short_id: A-SEPARATE
% verification: never
% category: protocol-structure
% affects: [secrecy, correctness]
% user_statement: >-
%   You run several executions, one after another or at the same time, and
%   want their security parameters to add. Their number and parameters are
%   fixed before the first starts. Each execution has its own key material,
%   randomness, seeds and pre-shared authentication key, such as a separate
%   part of one pool. No device, memory or random source carries information
%   between them. An execution has ended once both parties have fixed
%   their output, key or abort, after the closing messages; no key is used
%   or read, and nothing else leaks from an ended one, until all have ended.
%   The
%   assumptions each execution's run record lists must still hold when an
%   attacker holds and runs every other execution, her memory, which then
%   holds their systems, of any dimension. A key that
%   authenticates another execution, a key used early, or parameters chosen
%   from an earlier execution's data violate this.
% lean:
% ---

Composable frameworks add the errors of protocols composed in sequence or
in parallel \cite[Sec.~II.E]{portmann2022security}. For the trace-distance
parameter of this contract, the corresponding sum needs the executions to
be separate in the sense below, so that each can be analysed with every
other one on the adversary's side.

\begin{assumption}\label{ass:separate-executions}
For executions $1,\dots,N$ of the protocol, run against one adversary:
(i) $N$ and, for every execution, the parameters fixed before its run (its
acceptance set, round count, output length, transcript length, error
budgets and authentication parameters) are fixed before the first of the
$N$ executions starts;
(ii) no operation of the honest parties of one execution acts on a register
that the honest parties of another execution use, hold or output: no
execution uses any part of another execution's key, local randomness,
seeds or pre-shared key, no device, memory or random source carries a
state from one execution into another, and the executions interact only
through systems and classical channels the adversary has access to;
(iii) from the \tref{end of an execution}
(Definition~\ref{def:execution-end}) until the end of the last of the $N$
executions, no part of that execution's key is used, read or otherwise
acted on, and no information about any other register of its parties
leaves their laboratories;
(iv) for every execution, every assumption on which its key and its
security parameter rest holds when every register of the honest parties of
every other execution is counted as part of the adversary's information,
and every operation of those parties as part of the adversary's actions.
\end{assumption}
